\section{Self-exciting point processes for incident streams}\label{sec:framework}

This section fixes the model class used throughout the paper. We first recall the
likelihood theory of multivariate point processes. We then define a Hawkes process
augmented by exogenous shocks and recall its branching (cluster) representation and
its first- and second-order structure. The results in this section are classical; we
state them in the form needed later, with attributions and short proofs where the
argument is brief. \Cref{sec:th-meaning} explains what ``contagion'' means, and does
not mean, when the model is applied to antisemitic incidents.

\subsection{Marked multivariate point processes}\label{sec:th-mpp}

We observe $M$ \emph{streams} of incidents on a window $[0,T]$, indexed by $k\in[M]$.
A stream is any class of events that we want to model jointly. For example, offline
incidents may be split by region and by offence type (assault, vandalism, harassment,
threats). An online channel can also be a stream, with one point for each item that a
classifier flags. The points of stream $k$ are $0<t^k_1<t^k_2<\dots\le T$, its counting
process is $N_k(t)=\#\{i:t^k_i\le t\}$, and each point may carry a mark
$x^k_i\in\mathbb X$ (for example severity, target type or source). We assume that
every $N_k$ is simple and that no two streams share a point, almost surely. The history
$\Hist_t$ is the $\sigma$-field generated by all streams and marks on $[0,t]$, together
with the exogenous events $(\tau_e,S_e)$ that have occurred by time $t$ (defined below).

The \emph{conditional intensity} of stream $k$ is the nonnegative, $\Hist$-predictable
process $\lambda_k$ for which
\[
N_k(t)-\Lam_k(t),\qquad \Lam_k(t)=\int_0^t\lambda_k(s)\,\dd s,
\]
is a local martingale. Heuristically,
$\lambda_k(t)\,\dd t=\Pr\{\dN_k(t)=1\mid\Hist_{t-}\}$. The process $\Lam_k$ is the
\emph{compensator}. For a parametric family of intensities $\lambda_k(\cdot;\vartheta)$,
the log-likelihood of the observed streams on $[0,T]$ is
\begin{equation}\label{eq:th-loglik}
  \ell_T(\vartheta)=\sum_{k=1}^{M}\Big[\sum_{i:\,t^k_i\le T}\log\lambda_k(t^k_i;\vartheta)
  -\int_0^T\lambda_k(t;\vartheta)\,\dd t\Big]
\end{equation}
\citep[Ch.~7]{daley2003}; see \citet{ogata1978} for its asymptotic theory and
\citet{ogata1988} for its use in self-exciting models. If marks are
\emph{unpredictable}, meaning that the mark of a new point is drawn from a law that does
not depend on the past, the mark density adds a separate term to \eqref{eq:th-loglik}.
The compensator also drives the time-rescaling diagnostics of \cref{sec:estimation}.

\subsection{Hawkes processes with exogenous shocks}\label{sec:th-hawkes}

Salient external events, such as geopolitical escalations, terror attacks or
high-profile public statements, are indexed by $e$. Event $e$ occurs at time $\tau_e$
and has size $S_e\ge0$. Our working model is the linear multivariate Hawkes process
\citep{hawkes1971} augmented by these exogenous shocks:
\begin{equation}\label{eq:th-intensity}
  \lambda_k(t)=\mu_k(t)+\sum_{e:\,\tau_e<t}\upsilon_k S_e\,\psi_k(t-\tau_e)
  +\sum_{j=1}^{M}\int_{0}^{t-}\varphi_{kj}(t-s)\,\dN_j(s),
  \qquad \varphi_{kj}=\alpha_{kj}\,g_{kj}.
\end{equation}

\begin{assumption}[Model]\label{ass:th-model}
\begin{enumerate}[(i)]
  \item The background $\mu_k(t)=\exp(B(t)^{\top}\theta_k)$ is built from a spline or
    seasonal basis $B$ and is bounded between two positive constants on $[0,T]$.
    Calendar effects such as holiday periods belong to $\mu_k$, not to the shocks.
  \item The shock responses satisfy $\psi_k\ge0$, $\int\psi_k=1$ and $\psi_k(u)=0$
    for $u<0$. The loadings satisfy $\upsilon_k\ge0$.
  \item Each $g_{kj}$ is a probability density on $(0,\infty)$, and $\alpha_{kj}\ge0$.
\end{enumerate}
\end{assumption}

The constant $\alpha_{kj}$ is the expected number of stream-$k$ incidents directly
triggered by one stream-$j$ incident. These constants form the \emph{branching matrix}
$\Amat=[\alpha_{kj}]$, and the kernels form the matrix
$\Phik(t)=[\varphi_{kj}(t)]$. The \emph{resolvent} is
$\Rker=\sum_{m\ge1}\Phik^{*m}$, where $*$ denotes convolution. When the spectral radius
$\spr(\Amat)$ is below one, $\int\Rker=\sum_{m\ge1}\Amat^m$ is finite and the
\emph{contagion multiplier} is
\[
\Mmat=(\Imat-\Amat)^{-1}=\Imat+\int_0^\infty\Rker(t)\,\dd t=\sum_{m\ge0}\Amat^m .
\]
The workhorse special case uses exponential kernels,
$g_{kj}(t)=\beta_{kj}e^{-\beta_{kj}t}$. In the univariate case,
$\varphi(t)=n\beta e^{-\beta t}$, the branching ratio is $n=\alpha$ and, when $n<1$,
the stationary mean rate is $\bar\Lambda=\mu/(1-n)$. Throughout, $\beta$ (or
$\beta_{kj}$) is a decay rate and $\upsilon_k$ a shock loading.

\subsection{The cluster representation}\label{sec:th-cluster}

\Citet{hawkes1974} showed that a self-exciting process is a Poisson cluster process.
For \eqref{eq:th-intensity} the construction is as follows.
\begin{enumerate}[(a)]
  \item \emph{Immigrants.} Stream $k$ receives \emph{background immigrants}, a Poisson
    process with intensity $\mu_k(t)$. Independently, it receives \emph{shock
    immigrants}, a Poisson process with intensity
    $\sum_e\upsilon_kS_e\psi_k(t-\tau_e)$.
  \item \emph{Offspring.} Each point of stream $j$ at time $s$, whether an immigrant or
    not, independently generates offspring in each stream $k$. These offspring form a
    Poisson process with intensity $\varphi_{kj}(t-s)$, $t>s$. Equivalently, there are
    $\Pois(\alpha_{kj})$ offspring, placed at independent delays with
    density $g_{kj}$.
  \item \emph{Recursion.} Offspring generate offspring in the same way.
\end{enumerate}
The superposition of all generations has intensity \eqref{eq:th-intensity} with respect
to its own history and the shock calendar \citep{hawkes1974,reinhart2018,laub2021}.
Each incident belongs to exactly one cluster, the family tree of its immigrant
ancestor. The expected number of generation-$g$ descendants in stream $k$ of a
stream-$j$ incident is $(\Amat^g)_{kj}$. Hence $\Mmat_{kj}$ is the expected number of
stream-$k$ incidents in a cluster rooted at a stream-$j$ incident, the root included.
The latent parent labels drive the EM algorithm of \cref{sec:estimation}. The split of
incidents into background immigrants, shock immigrants and offspring underlies the
surge decomposition of \cref{sec:attribution}.

\begin{proposition}[Stationarity; \citealp{hawkes1974,bremaud1996}]\label{prop:stationary}
Suppose there are no shocks, $\mu_k(t)\equiv\mu_k>0$ for every $k$, and
\cref{ass:th-model}(iii) holds.
\begin{enumerate}[(a)]
  \item If $\spr(\Amat)<1$, there is a stationary point process on $\mathbb R$ whose
    intensity with respect to its own history is
    $\lambda_k(t)=\mu_k+\sum_j\int_{(-\infty,t)}\varphi_{kj}(t-s)\,\dN_j(s)$. Its law is
    unique among stationary solutions with finite mean intensity. It is the Poisson
    cluster process of (a)--(c), and its clusters are almost surely finite. Its mean
    intensity vector is
    \[
    \bar{\bm\Lambda}=(\Imat-\Amat)^{-1}\bm\mu=\Mmat\bm\mu .
    \]
    The expected size of a cluster rooted at a stream-$j$ immigrant is
    $\mathbf 1^{\top}\Mmat\bm e_j$. The long-run fraction of stream-$k$ incidents that
    are offspring is $1-\mu_k/\bar\Lambda_k=(\Amat\bar{\bm\Lambda})_k/\bar\Lambda_k$.
  \item Conversely, if a stationary solution with finite mean intensity exists, then
    $\spr(\Amat)<1$.
  \item In the univariate case, the cluster size $\mathcal S$ (immigrant included) has the
    Borel distribution
    $\Pr(\mathcal S=k)=e^{-nk}(nk)^{k-1}/k!$, $k\ge1$, with $\E \mathcal S=(1-n)^{-1}$,
    $\E \mathcal S^2=(1-n)^{-3}$, $\E \mathcal S^3=(1+2n)(1-n)^{-5}$ and
    $\E \mathcal S^4=(1+8n+6n^2)(1-n)^{-7}$. Moreover $\E e^{u \mathcal S}<\infty$ for
    $u\le n-1-\log n$, and $n-1-\log n>0$ for $n\in(0,1)$.
\end{enumerate}
\end{proposition}

The proof is in \cref{app:fw-proofs}.

Part (b) uses $\mu_k>0$. Without immigrants, critical processes with finite mean
intensity can be stationary under heavy-tailed kernels \citep{bremaud2001}. In our
setting $\mu_k>0$ excludes this case.

\subsection{Second-order structure}\label{sec:th-second}

Let $\bm N$ be a stationary $M$-variate point process with mean vector
$\bar{\bm\Lambda}$ and finite second moments. Its covariance structure has two parts.
The first is an atom $\operatorname{diag}(\bar{\bm\Lambda})\,\delta_0$ at lag zero,
which comes from each point paired with itself. The second is, we assume, a covariance
density $\bm c(u)=[c_{kj}(u)]$:
\[
\Cov\{\dN_k(t+u),\dN_j(t)\}=\big[\delta_{kj}\bar\Lambda_k\,\delta_0(u)+c_{kj}(u)\big]\,\dd u\,\dd t .
\]
The \emph{Bartlett spectrum} \citep{bartlett1963} is the Fourier transform of this
covariance measure:
\begin{equation}\label{eq:th-bartlett-def}
\bm f(\omega)=\frac1{2\pi}\Big[\operatorname{diag}(\bar{\bm\Lambda})+\hat{\bm c}(\omega)\Big],
\qquad \hat{\bm c}(\omega)=\int_{\mathbb R}e^{-i\omega u}\bm c(u)\,\dd u
\end{equation}
\citep[Ch.~8]{daley2003}. A Poisson process has the flat spectrum
$\bar\Lambda/(2\pi)$. Every second-order statistic of counts is a functional of the
pair $(\bar{\bm\Lambda},\bm f)$. This includes the autocovariances of the binned
counts $X_t=N((t-1)\Delta,t\Delta]$ for any bin width $\Delta$, the variance-time curve
$w\mapsto\Var N(0,w]$, the expected periodogram, and the asymptotic index of
dispersion
\[
F_\infty=\lim_{w\to\infty}\frac{\Var N(0,w]}{\E N(0,w]}=\frac{2\pi f(0)}{\bar\Lambda}.
\]

\begin{proposition}[Bartlett spectrum of a Hawkes process; \citealp{hawkes1971}]\label{prop:th-spectrum}
Under the conditions of \cref{prop:stationary}(a), the stationary Hawkes process has
finite moments of all orders, and its Bartlett spectrum is
\begin{equation}\label{eq:th-bartlett}
  \bm f(\omega)=\frac1{2\pi}\big(\Imat-\hat\Phik(\omega)\big)^{-1}
  \operatorname{diag}(\bar{\bm\Lambda})\,\big(\Imat-\hat\Phik(\omega)\big)^{-*},
  \qquad \hat\Phik(\omega)=\int_0^\infty e^{-i\omega t}\Phik(t)\,\dd t ,
\end{equation}
where $\bm B^{-*}=(\bm B^{-1})^{*}$ denotes the inverse of the conjugate transpose. In
the univariate case, the covariance density is
$c(u)=\bar\Lambda\big[R(|u|)+\int_0^\infty R(v)R(v+|u|)\,\dd v\big]$ and
$F_\infty=(1-n)^{-2}$, whatever the kernel. For the exponential kernel,
\begin{equation}\label{eq:th-fH}
  f_H(\omega)=\frac{\bar\Lambda}{2\pi}\Big[1+\frac{\beta^2n(2-n)}{\beta^2(1-n)^2+\omega^2}\Big],
  \qquad
  c(u)=\frac{\bar\Lambda\,\beta\,n(2-n)}{2(1-n)}\,e^{-\beta(1-n)|u|}.
\end{equation}
Here $n(2-n)=1-(1-n)^2$.
\end{proposition}

A proof sketch is in \cref{app:fw-proofs}.

Two features of \eqref{eq:th-fH} matter later. First, the correlation time of the
counts is $1/\{\beta(1-n)\}$, which is longer than the kernel time scale $1/\beta$:
cascades lengthen the memory. Second, the excess dispersion is
$F_\infty-1=n(2-n)/(1-n)^2$, which grows without bound as $n\uparrow1$.

\subsection{Interpretation for antisemitic incidents}\label{sec:th-meaning}

\paragraph{Streams.} Offline streams may be incidents recorded by police or monitoring
organisations, split by region and by offence type. An online stream may be the volume
of items that a classifier labels antisemitic on one platform. It can be modelled
either as a point process or, after binning, through the discrete-time approximation
of \cref{sec:estimation} \citep{kirchner2016}. Online volume is a \emph{measured}
quantity whose errors matter for inference (\cref{rem:th-granger-caveats}).

\paragraph{What self-excitation can represent.} A positive $\alpha_{kj}$ says that,
under the model, a stream-$j$ incident raises the short-run probability of further
stream-$k$ incidents. Several mechanisms could produce this, and each is a
\emph{hypothesis} that the model does not by itself verify:
\begin{itemize}
  \item \emph{imitation}, or copycat behaviour;
  \item \emph{emboldening}, where visible incidents shift perceived norms about what is
    tolerated;
  \item \emph{network diffusion} through online communities, coordinated campaigns or
    repeat offenders;
  \item \emph{repeat or near-repeat targeting} of the same location, as in
    self-exciting models of crime \citep{mohler2011}.
\end{itemize}
Cross-excitation from an online to an offline stream is the formal counterpart of the
hypothesis that online hostility precedes offline incidents.

\paragraph{What contagion does not mean.} Three limits apply.
\begin{itemize}
  \item The parameters describe predictive dependence in a probability model.
    ``Contagion'' is used in its technical sense of self-excitation; compare the
    classical distinction between true and apparent contagion \citep{feller1943}.
    It attributes no moral responsibility to victims, to Jewish communities, or to anyone
    connected with an earlier incident. Responsibility for an act rests with its
    perpetrator.
  \item A positive $\alpha_{kj}$ is not proof of a causal mechanism.
    \Cref{sec:identification} shows that, within a single stream, unobserved shocks can
    mimic contagion exactly at second order, for exponential and completely monotone
    kernels. It also shows that under-recording attenuates contagion,
    and that common drivers can induce online-to-offline Granger causality without any
    direct link.
  \item Shocks are modelled as \emph{salient external events} whose timing is taken
    from an external calendar. Including an event in the calendar, and estimating its
    loading, is a statement about incident dynamics. It is not a judgement about the
    event, its actors or its merits.
\end{itemize}
